\section{Arquitecturas de redes neuronales convolucionales (CNN)}

\subsection{Los tres componentes centrales de una CNN}

Habiendo entendido la operación de convolución, podemos construir una red neuronal convolucional completa. Una CNN típica está compuesta por tres tipos de capas que se alternan:

\begin{figure}[H]
\centering
\includegraphics[width=0.9\textwidth]{slides-images/slide-44.jpg}
\caption{Flujo completo de una red CNN de clasificación: imagen de entrada $\rightarrow$ capa convolucional (genera mapas de características) $\rightarrow$ max pooling $\rightarrow$ $\cdots$ $\rightarrow$ capa totalmente conectada $\rightarrow$ salida de clasificación.\protect\footnotemark}
\end{figure}
\footnotetext{Diapositivas, página 44.}

\begin{importantbox}{Las tres operaciones centrales de una CNN}
\begin{enumerate}
    \item \textbf{Convolución} (Convolution): aplica un filtro para generar un mapa de características, extrayendo características locales
    \item \textbf{Activación no lineal} (Non-linearity): normalmente se usa ReLU, que introduce capacidad de representación no lineal
    \item \textbf{Pooling} (Pooling): submuestrea cada mapa de características, reduciendo la dimensión y obteniendo invariancia espacial
\end{enumerate}
Mediante el entrenamiento, la CNN aprende automáticamente los pesos de los filtros de las capas convolucionales.
\end{importantbox}

\subsection{La capa convolucional: conectividad local}

Cada neurona de una capa convolucional se conecta únicamente a una región local de la entrada. Para la neurona en la posición $(p,q)$ de la capa oculta:

\begin{figure}[H]
\centering
\includegraphics[width=0.85\textwidth]{slides-images/slide-46.jpg}
\caption{Conectividad local de la capa convolucional: cada neurona toma su entrada de una región local, calcula la suma ponderada y agrega el sesgo. Su expresión matemática es la misma que la de una capa totalmente conectada, pero su alcance es local.\protect\footnotemark}
\end{figure}
\footnotetext{Diapositivas, página 46.}

$$y_{p,q} = \sum_{i=1}^{k} \sum_{j=1}^{k} w_{ij} \cdot x_{i+p, j+q} + b$$

Este proceso consta de tres pasos: (1) aplicar la ventana de pesos, (2) calcular la combinación lineal, (3) activar mediante una función no lineal.

\subsection{Disposición espacial del volumen de salida}

La salida de cada capa no es un plano bidimensional, sino un \textbf{volumen} (volume) tridimensional, de dimensiones $h \times w \times d$:

\begin{figure}[H]
\centering
\includegraphics[width=0.9\textwidth]{slides-images/slide-47.jpg}
\caption{Dimensiones del volumen de salida de una CNN: $h \times w \times d$, donde $h$ y $w$ son las dimensiones espaciales y $d$ (la profundidad) es igual al número de filtros de esa capa.\protect\footnotemark}
\end{figure}
\footnotetext{Diapositivas, página 47.}

\begin{itemize}
    \item $h$ y $w$: las dimensiones espaciales de la salida (determinadas por el tamaño de la entrada, el tamaño del filtro y el stride)
    \item $d$: la profundidad, igual al \textbf{número de filtros} usado en esa capa
    \item \textbf{Stride} (paso): el tamaño del desplazamiento del filtro al deslizarse
    \item \textbf{Receptive Field} (campo receptivo): la región de la imagen de entrada conectada a un nodo de salida determinado
\end{itemize}

En TensorFlow: \texttt{tf.keras.layers.Conv2D(filters=d, kernel\_size=(h,w), strides=s)}

En PyTorch: \texttt{torch.nn.Conv2d(in\_channels, out\_channels=d, kernel\_size, stride=s)}

\begin{warningbox}{PyTorch requiere especificar el número de canales de entrada}
El \texttt{Conv2D} de TensorFlow/Keras infiere automáticamente la dimensión de entrada, pero el \texttt{Conv2d} de PyTorch exige especificar explícitamente \texttt{in\_channels}. Quienes empiezan suelen olvidar esto y obtienen un error.
\end{warningbox}

\subsection{Función de activación no lineal: ReLU}

Después de cada operación de convolución es necesario aplicar una función de activación no lineal. En las CNN, la más usada es \textbf{ReLU} (Rectified Linear Unit):

\begin{figure}[H]
\centering
\includegraphics[width=0.85\textwidth]{slides-images/slide-48.jpg}
\caption{Función de activación ReLU: reemplaza todos los valores negativos por cero y conserva los valores positivos sin cambio. Es una operación no lineal aplicada píxel por píxel.\protect\footnotemark}
\end{figure}
\footnotetext{Diapositivas, página 48.}

$$g(z) = \max(0, z)$$

La función de ReLU consiste en reemplazar por cero todos los valores negativos del mapa de características, conservando solo los valores positivos (es decir, las regiones donde se detectó la característica). ¿Por qué ReLU resulta especialmente adecuada para tareas con imágenes? Porque cuando un filtro detecta un patrón determinado, la salida es positiva (coincidencia), y cuando no lo detecta, la salida es negativa o cero (no hay coincidencia); ReLU conserva precisamente la señal de «se detectó».

\subsection{Pooling: reducción de dimensión e invariancia espacial}

El \textbf{pooling} (Pooling) es una operación de submuestreo que sigue a la capa convolucional. La más usada es el \textbf{max pooling} (Max Pooling):

\begin{figure}[H]
\centering
\includegraphics[width=0.85\textwidth]{slides-images/slide-49.jpg}
\caption{Ejemplo de max pooling: usando una ventana de 2$\times$2 y un stride de 2, se toma el valor máximo en cada ventana. Una entrada de 4$\times$4 se convierte en una salida de 2$\times$2.\protect\footnotemark}
\end{figure}
\footnotetext{Diapositivas, página 49.}

\begin{importantbox}{Las dos funciones principales del pooling}
\begin{enumerate}
    \item \textbf{Reducción de dimensión} (Reduced dimensionality): reduce el tamaño espacial del mapa de características, disminuyendo así el costo de cómputo y el número de parámetros de las capas siguientes
    \item \textbf{Invariancia espacial} (Spatial invariance): al tomar el valor máximo, la posición exacta de la característica importa menos, lo que aumenta la robustez frente a desplazamientos
\end{enumerate}
\end{importantbox}

Amini menciona que, además del max pooling, se puede plantear una pregunta más general: «¿de qué otras formas se puede lograr el submuestreo y a la vez mantener la invariancia espacial?». Esto sugiere otras técnicas que quizá se presenten en clases posteriores, como la convolución con stride (strided convolution).

\subsection{Aprendizaje de representaciones en CNN profundas}

Al apilar varios módulos de «convolución + ReLU + pooling», la CNN es capaz de aprender representaciones de características cada vez más abstractas:

\begin{figure}[H]
\centering
\includegraphics[width=0.85\textwidth]{slides-images/slide-50.jpg}
\caption{Aprendizaje de representaciones en una CNN profunda: la primera capa aprende características de bajo nivel como bordes y manchas oscuras, la segunda las combina en características de nivel medio como ojos y nariz, y la tercera las combina en características de alto nivel como la estructura facial completa.\protect\footnotemark}
\end{figure}
\footnotetext{Diapositivas, página 50. Referencia: Lee y otros, ICML 2009.}

\begin{knowledgebox}{Composición jerárquica de características}
Las CNN profundas son poderosas porque cada capa realiza una \textbf{composición} sobre las características detectadas por la capa anterior:
\begin{itemize}
    \item La primera capa convolucional detecta características primitivas (como bordes en distintas direcciones)
    \item Después del pooling, el mapa de características se reduce, lo que equivale a observar a una escala espacial mayor
    \item La segunda capa convolucional detecta sobre el mapa de características reducido, lo que equivale a combinar varias características de bajo nivel en características de nivel medio
    \item Al repetirse este proceso, las características se vuelven cada vez más abstractas y semánticamente más ricas
\end{itemize}
Este aprendizaje jerárquico de características, de lo simple a lo complejo, es precisamente lo que significa la palabra «profundo» en aprendizaje profundo.
\end{knowledgebox}

\subsection{Resumen de la sección}

Una red CNN de clasificación completa está compuesta por una parte de aprendizaje de características y una parte de clasificación. La parte de aprendizaje de características, alternando capas convolucionales (que extraen características locales), activación no lineal (que aporta capacidad de representación) y capas de pooling (que reducen la dimensión y aumentan la invariancia), construye capa por capa representaciones de características cada vez más abstractas. La parte de clasificación aplana las características aprendidas y las pasa por una capa totalmente conectada para obtener las probabilidades de clase.

\newpage
